% Texto: borrador de Laura Segura, con acentos restaurados y citas enlazadas.

\section{Materiales y métodos}

\subsection{Tipo de investigación y enfoque}

La investigación fue de tipo aplicada y de diseño tecnológico. El objeto
evaluado no fue la optimización predictiva del modelo, sino la capacidad del
marco para generar evidencia verificable de cumplimiento. Por ello, el modelo de
aprendizaje automático se trató como sujeto de prueba: se entrenó y selló antes
de ejecutar el procedimiento de cumplimiento, y no se ajustó posteriormente con
base en los resultados de equidad o explicabilidad.

La metodología se organizó en cuatro componentes. El primero fue una matriz de
mapeo principio-requerimiento. El segundo correspondió a tres módulos de
cumplimiento: explicabilidad, equidad y trazabilidad. El tercero fue un registro
estructurado de inferencias. El cuarto fue un procedimiento reproducible de seis
pasos para generar el expediente de evidencia.

\subsection{Datos y población de estudio}

Se utilizaron datos públicos del proyecto CASAS, compuesto por eventos de
sensores ambientales en entornos domiciliarios \cite{cook2012aruba}. La primera
configuración corresponde al hogar Aruba, con una vivienda, sensores PIR,
sensores de puerta, sensores de temperatura y actividades anotadas. La segunda
corresponde a una serie multi-hogar, compuesta por nueve viviendas de personas
adultas mayores de una comunidad de retiro.
% El depósito de Zenodo de la serie hh pide citar Cook et al. (2013)
% [cook2013casas] y el propio conjunto [cook2025zenodo]; están en el .bib.
\pendiente{citar el depósito de la serie multi-hogar (Zenodo) y el artículo que pide citar}
En ambos casos, la población observada corresponde a personas adultas mayores
en viviendas independientes, aunque el proyecto no recolectó datos nuevos ni
tuvo contacto con participantes humanos.
% config.yaml describe a la residente de Aruba como "mujer adulta voluntaria":
% no dice que sea adulta mayor. Para la serie hh sí consta (comunidad de retiro).
\revisar{para Aruba, la documentación del conjunto dice ``mujer adulta'', no adulta mayor}

Los eventos se transformaron en ventanas disjuntas de 30 eventos consecutivos.
Cada ventana generó una fila tabular con conteos por zona del hogar, eventos de
puerta, variables temporales y, cuando estaban disponibles, lecturas de
temperatura. La partición de entrenamiento y prueba fue temporal por día, con
una proporción de prueba de 0,2. Esta decisión evitó mezclar ventanas contiguas
de una misma secuencia entre entrenamiento y prueba.

\subsection{Modelo de referencia}

El modelo de referencia fue un clasificador \texttt{RandomForestClassifier} con
semilla fija, ejecución de un solo hilo y versiones de bibliotecas declaradas.
Su función fue proporcionar inferencias para auditar el marco. La calidad
predictiva del modelo no se usó como criterio de éxito del estudio, pues el
propósito fue comprobar si el procedimiento producía evidencia completa y
reconstruible.
% Disponibles en el .bib: breiman2001rf (bosques aleatorios) y
% pedregosa2011sklearn (scikit-learn).

\subsection{Matriz de requerimientos}

La matriz tradujo los principios seleccionados en nueve requerimientos
operativos. El Cuadro~\ref{cuadro:requerimientos} resume la correspondencia
entre principio, requerimiento y evidencia esperada.

\begin{table}[htbp]
  \caption{Requerimientos operativos y evidencia del marco.}
  \label{cuadro:requerimientos}
  \small
  \begin{tabularx}{\linewidth}{@{}>{\raggedright\arraybackslash}p{0.27\linewidth}c>{\raggedright\arraybackslash}X@{}}
    \toprule
    Principio & Código & Evidencia técnica y documental \\
    \midrule
    Transparencia y explicabilidad & R3.1 & Explicación local por inferencia mediante atribución cuantificada de variables. \\
    Transparencia y explicabilidad & R3.2 & Importancia global agregada sobre el conjunto de evaluación. \\
    Transparencia y explicabilidad & R3.3 & Traducción de resultados a lenguaje comprensible y model card. \\
    \midrule
    Equidad y no discriminación & R4.1 & Desempeño desagregado por subgrupo. \\
    Equidad y no discriminación & R4.2 & Métricas de disparidad con umbrales declarados antes de medir. \\
    Equidad y no discriminación & R4.3 & Datasheet del conjunto de datos con procedencia, composición y limitaciones. \\
    \midrule
    Responsabilidad & R5.1 & Bitácora estructurada y persistente de inferencias. \\
    Responsabilidad & R5.2 & Model card con propósito, desempeño, limitaciones y condiciones de uso. \\
    Responsabilidad & R5.3 & Verificación de reconstrucción y atribución ex post de decisiones. \\
    \bottomrule
  \end{tabularx}
  \fuente{elaboración propia.}
\end{table}

\subsection{Módulo de explicabilidad}

El módulo de explicabilidad utilizó SHAP con TreeExplainer para calcular
atribuciones locales de las variables de entrada
\cite{lundberg2017shap,lundberg2020treeshap}. Para cada inferencia del conjunto
de prueba se registró la clase predicha, el valor base, las contribuciones de
las variables y una referencia al artefacto de explicación. La importancia
global se obtuvo como la media de las contribuciones absolutas locales, evitando
un segundo cálculo independiente y manteniendo coherencia entre explicaciones
locales y globales. Las variables se nombraron por zonas del hogar, por ejemplo
cocina o dormitorio, para favorecer la comprensión por parte de destinatarios no
técnicos.

\subsection{Módulo de equidad}

El módulo de equidad evaluó desempeño desagregado y disparidad entre subgrupos.
En Aruba, los subgrupos fueron contextuales, como franja horaria y tipo de día.
En la configuración multi-hogar, se incorporó el hogar como subgrupo
poblacional, pues cada vivienda representa una persona o unidad doméstica
distinta. Las métricas con efecto en el veredicto fueron diferencia de tasa de
verdaderos positivos y diferencia de tasa de falsos positivos, ambas con umbral
de 0,10 declarado antes de la ejecución. También se reportaron como descriptivas
la diferencia de paridad demográfica y el cociente de tasa de selección.
% Disponible en el .bib: hardt2016equality, origen de la igualdad de
% oportunidades (diferencia de TVP).

\subsection{Módulo de trazabilidad}

La trazabilidad se implementó mediante una bitácora JSONL \emph{append-only}.
Cada línea registró identificador de evento, marca temporal UTC, hash de la
entrada procesada, salida del modelo, confianza, versión del modelo, versión del
marco, referencia a la explicación local, responsable declarado y versión del
esquema. El uso de hashes evitó persistir datos crudos sensibles en la
bitácora, sin impedir la reconstrucción técnica de la decisión dentro del
expediente.

\subsection{Procedimiento de aplicación}

El procedimiento siguió seis pasos alineados con las funciones del NIST AI RMF:
caracterización del sistema, documentación del conjunto de datos, declaración de
criterios de evaluación, ejecución de pruebas técnicas, generación de artefactos
auditables y verificación de auditabilidad. El paso de declaración de criterios
selló el hash SHA-256 del archivo de configuración. En la verificación final, el
marco comprobó que ese hash no hubiera cambiado, que la bitácora fuera válida,
que existiera cobertura uno a uno entre inferencias esperadas y registradas, y
que los nueve requerimientos contaran con artefactos presentes.
% Disponibles en el .bib: gebru2021datasheets y mitchell2019modelcards, que
% originan el datasheet y el model card del procedimiento.
